%% kapa: exam=1 round=32 year=2021 period=2 form=A subject=accounting subject_ko=회계학 seq=6 number=46
%% source: Q-Net 감정평가사 기출문제 (cst00309) artlSeq=5208397 file=extracted/1008/5208397/A형/2021년도 제32회 감정평가사 1차 2교시 A형.pdf
%% conversion: kordoc 4.24.51 -> tools/gen1.py (auto); review=proofread-sonnet; flags=none
\begin{question}{46}
(주)감평은 20×1년 초 해지불능 리스계약을 체결하고 사용권자산(내용연수 5년, 잔존가치 ￦0, 정액법 상각)과 리스부채(리스기간 5년, 매년 말 정기리스료 ￦13,870, 리스기간 종료 후 소유권 무상이전 약정)를 각각 ￦50,000씩 인식하였다. 리스계약의 내재이자율은 연 12\%이고 (주)감평은 리스회사의 내재이자율을 알고 있다. (주)감평은 사용권자산에 대해 재평가모형을 적용하고 있으며 20×1년 말 사용권자산의 공정가치는 ￦35,000이다. 동 리스계약이 (주)감평의 20×1년 당기순이익에 미치는 영향은? (단, 리스계약은 소액자산리스 및 단기리스가 아니라고 가정한다.)
\begin{choices}
  \item ￦5,000 감소
  \item ￦6,000 감소
  \item ￦15,000 감소
  \item ￦16,000 감소
  \item ￦21,000 감소
\end{choices}
\end{question}
\answer{5}
